\chapter{Classification}\label{sec-classification}

The term `class' originates from the Latin \textit{classis}, which was used in
ancient Rome to categorize citizens within the Centuriate Assembly based
on wealth and property. This system, attributed to the reforms of
Servius Tullius, the sixth king of Rome, divided society into several
classes. The wealthiest in the upper classes received payments to equip themselves as members of the cavalry and were granted significant political power. In contrast, the lower
classes consisted of those with less wealth, who had to arm themselves and with limited political influence. This classification
primarily determined military duties and voting rights, reflecting the
organization of society by economic status. Later, this classification
of individual citizens was used for taxation. The assignment of
a citizen to one of the five existing classes took into account several factors,
such as the type or armor they possessed, wealth, and age, among others.
Classification for taxation rather than for military reasons was important
because there was a constant need for revenue, whereas wars occurred less frequently (Cornell 2012).

Some two thousand five hundred years later we have developed
several methods of of classification, one of them being
manual. Due to the incredible amount of data we deal with, however, this is very
impractical for the processing of information. In Chapter 5 we explored a method
known as clustering. Both classification and clustering aim to achieve
the same goal, that of putting data into categories using algorithms. Clustering
uses unsupervised learning, where the algorithm, with the help of
mathematical formulas, attempts to group elements with common
characteristics without predefined categories. This means we interpret
the categories after the algorithm has done its work.

On the other hand, classification involves supervised learning. Here, we
start with a training set, a collection of documents already categorized
into classes or categories. The algorithm learns from these existing
categories to create a model, or a `modus operandi,' through which to categorize
the new data. This process is essential when data needs to be sorted
automatically to save time and achieve accurate results.

For instance, a customer-relationship manager dealing with hundreds of
incoming emails daily could use an automated system to sort these emails
into the correct departments, such as: Finance, Logistics, or Complaints. Another classic application in a direct marketing
context is RFM classification, which involves segmenting customers based
on the recency, the frequency, and monetary value of their previous
interactions with a firm (see, e.g., Bult and Wansbeek 1995; Gönül, Kim, and Shi 2000; Berry and Linoff 2009). The application
of automatic classification is many and varied.

This chapter will focus on classification using the same tablet data set
used in the previous chapters. Note that there is a `Brand' column at
the end of the dataset, which we will use for classification. Each
review in the data set is associated with a brand. The goal is for every
new review which lacks a classification by brand, to be processed by the
model and assigned to the correct category of brand.

This chapter will
cover the following types of classification algorithms: linear discriminant analysis, K-nearest neighbors, support vector
machines, and neural networks. To complete the chapter, we will use a
library, FastText, for the efficient learning of word representation
and sentence classification. It is optimal for its speed of execution
and accuracy, especially with large datasets.

\section{Supervised machine-learning-algorithm
models}\label{supervised-machine-learning-algorithm-models}

In this section, we will focus on employing supervised machine-learning
algorithms. Our primary elements are the text content of the review and
its association with a specific brand. Using these well-understood
elements, the algorithm will generate a model. This enables the
classification of new, unlabeled content, ensuring that the correct
brand label is assigned accurately.

The entire process is depicted in Figure~\ref{fig-ml}, showing a
typical machine-learning workflow when applied to a corpus of documents. This
procedure is consistent across different types of data sets.
Initially (step I in Figure~\ref{fig-ml}), we provide the algorithm with
a set of documents which have already been classified, each review tagged
with a brand label. This data set is known as the training set,
offering the algorithm a comprehensive framework for how reviews
correlate with brands. Through the analysis of these associations and
specific rules governing classification, the machine constructs a model based on
the data, particularly the words within each review.

In the subsequent step, II, we introduce a collection of unlabeled
documents to the model already established. The expectation is for
these documents to be accurately classified according to the model's
guidelines. Computers excel when identifying regularities within data,
using these patterns to make efficient and precise predictions.
Essentially, this is a task of predicting: Given a set of reviews, the
machine will, using a certain degree of probability, predict the
appropriate assignment of the review to the brand.

The algorithm's accuracy is evaluated using a testing set, typically
comprising 20\% of the total data set. As shown in Figure~\ref{fig-ml},
this evaluation is an integral part of the workflow. After the training
phase (step I), the classification model is applied to the testing set
(step II), which is processed by the model as if it were unlabeled. The
outcomes of this classification are then compared against the actual
labels (known as ground-truth values). While the train-test split is commonly
80:20, other configurations are also possible, such as a training,
validation, and test split (e.g., 80:10:10). In these setups, the
validation set is used to fine-tune the model, while the test set is
reserved exclusively for evaluating its performance. The higher the
congruence between how algorithm assigns the material and the original labels,
the greater the algorithm's accuracy.

\begin{figure}

\centering{

\includegraphics[width=0.9\textwidth,height=\textheight]{figure/chapter6/ml2.png}

}

\caption{\label{fig-ml}Supervised Learning}

\end{figure}%

With the current state of technology it is not possible to guarantee that algorithms will perform flawlessly. However, they are efficient enough to
reduce manual effort significantly. When data are misclassified, human intervention
becomes necessary. A robust approach includes
regularly updating the training data set after correcting any errors.
This process not only introduces new vocabulary and scenarios to the
model but also enhances the accuracy of its predictions.

Classification algorithms vary in effectiveness, with some outperforming
others based on the specific requirements of the data and the context of
how it is applied. Broadly, we can categorize these algorithms into two
types: linear and nonlinear. Linear classification uses a straight line
(or a hyperplane in higher dimensions) to divide data into distinct
classes, which, obviously, works well for linearly separable data. On the other
hand, nonlinear classification employs more complex boundaries to
separate classes, enabling it to handle data with intricate patterns or
overlaps.

\subsection{Data-set preparation}\label{dataset-preparation}

As discussed in previous chapters, text data must be converted into a
vectorized format and organized into a matrix or, more simply, a table.
This format allows the computer to process the numerical data
efficiently. The process begins by importing the data set into the R
programming environment from a local source or the internet. (If you
needs to re-familiarize yourself with this step, re-read the introductory sections
of this book for detailed guidance.) Next, we focus on a subset of the
data by selecting two brands for analysis, such as Apple and Samsung,
though you could choose any brands featured in the data set which are
of interest. This selection effectively narrows down the data set
horizontally. We then refine the data set further by retaining only the
columns essential for classification, namely, `Content' and
`Brand'. This involves slicing the data set vertically, ensuring
we have a focused and manageable data set for our analysis. But first we need to load the necessary libraries and the data set.

\begin{Shaded}
\begin{Highlighting}[]
\CommentTok{\# Set global option to avoid automatic conversion of strings to factors}
\FunctionTok{options}\NormalTok{(}\AttributeTok{stringsAsFactors =} \ConstantTok{FALSE}\NormalTok{)}

\CommentTok{\# Load necessary libraries}
\FunctionTok{library}\NormalTok{(tm)}
\FunctionTok{library}\NormalTok{(caret)}
\FunctionTok{library}\NormalTok{(tidyverse)}
\FunctionTok{library}\NormalTok{(SnowballC)}
\FunctionTok{library}\NormalTok{(textstem)}
\FunctionTok{library}\NormalTok{(kableExtra)}
\FunctionTok{library}\NormalTok{(caretEnsemble)}

\CommentTok{\# Read the classified dataset, directly filter, and select required data}
\NormalTok{dt }\OtherTok{\textless{}{-}} \FunctionTok{read.csv}\NormalTok{(}\StringTok{"data/Product Reviews.csv"}\NormalTok{) }\SpecialCharTok{\%\textgreater{}\%}
  \FunctionTok{filter}\NormalTok{(Brand }\SpecialCharTok{\%in\%} \FunctionTok{c}\NormalTok{(}\StringTok{"Apple"}\NormalTok{, }\StringTok{"Samsung"}\NormalTok{)) }\SpecialCharTok{\%\textgreater{}\%}
  \FunctionTok{select}\NormalTok{(Content, Brand)}
\end{Highlighting}
\end{Shaded}

With a data set comprising two columns, one of which identifies two
brands, we have established a division: One segment will serve as the training
set, and the other as the test set. The test set allows us to simulate a
scenario where the test portion is unlabeled to evaluate the algorithm's
performance. By applying the model to the test set, we can compare its
predictions against the known classifications. By now, preparing the
dataset should be familiar to you and certain commands are provided for you to explore during the training sessions. Using
the \texttt{\%\textgreater{}\%} operator, pre-processing and other
related tasks are integrated seamlessly within a document term
matrix, one for the training set and another for the test
set. We encourage you to read the comments embedded in the code to
understand each step better. You can also comment and remove the comments from the
functions and check the results.

\begin{Shaded}
\begin{Highlighting}[]
\CommentTok{\# Set the seed for reproducibility}
\NormalTok{seed }\OtherTok{\textless{}{-}} \DecValTok{123}
\FunctionTok{set.seed}\NormalTok{(seed)}

\CommentTok{\# Assign "train" or "test" to each observation in the dataset}
\NormalTok{dt}\SpecialCharTok{$}\NormalTok{Type }\OtherTok{\textless{}{-}} \FunctionTok{sample}\NormalTok{(}\FunctionTok{c}\NormalTok{(}\StringTok{"train"}\NormalTok{, }\StringTok{"test"}\NormalTok{), }\FunctionTok{nrow}\NormalTok{(dt), }\AttributeTok{replace =} \ConstantTok{TRUE}\NormalTok{, }\AttributeTok{prob =} \FunctionTok{c}\NormalTok{(}\FloatTok{0.8}\NormalTok{, }\FloatTok{0.2}\NormalTok{))}

\CommentTok{\# Check the distribution of the brands over the Type variable}
\FunctionTok{table}\NormalTok{(dt}\SpecialCharTok{$}\NormalTok{Brand, dt}\SpecialCharTok{$}\NormalTok{Type)}
\end{Highlighting}
\end{Shaded}

\begin{verbatim}
         
          test train
  Apple    372  1520
  Samsung  365  1419
\end{verbatim}

We see notice the distribution of the brands, with the number of reviews
for both brands being fairly similar.

\begin{Shaded}
\begin{Highlighting}[]
\CommentTok{\# Create a corpus from the Content column for preprocessing}
\FunctionTok{library}\NormalTok{(tm)}
\FunctionTok{library}\NormalTok{(SnowballC)}
\NormalTok{corpus }\OtherTok{\textless{}{-}} \FunctionTok{VCorpus}\NormalTok{(}\FunctionTok{VectorSource}\NormalTok{(dt}\SpecialCharTok{$}\NormalTok{Content))}

\CommentTok{\# Preprocess the corpus and create a Document Term Matrix}
\NormalTok{dtm }\OtherTok{\textless{}{-}}\NormalTok{ corpus }\SpecialCharTok{\%\textgreater{}\%}
  \FunctionTok{tm\_map}\NormalTok{(}\FunctionTok{content\_transformer}\NormalTok{(tolower)) }\SpecialCharTok{\%\textgreater{}\%} \CommentTok{\# Convert to lowercase}
  \FunctionTok{tm\_map}\NormalTok{(}\FunctionTok{content\_transformer}\NormalTok{(removeWords), }\FunctionTok{stopwords}\NormalTok{(}\StringTok{"english"}\NormalTok{)) }\SpecialCharTok{\%\textgreater{}\%} \CommentTok{\# Remove stopwords}
  \FunctionTok{tm\_map}\NormalTok{(removePunctuation) }\SpecialCharTok{\%\textgreater{}\%} \CommentTok{\# Remove punctuation}
  \FunctionTok{tm\_map}\NormalTok{(removeNumbers) }\SpecialCharTok{\%\textgreater{}\%} \CommentTok{\# Remove numbers}
  \CommentTok{\# tm\_map(stemDocument) \%\textgreater{}\% \# Stemming (commented out)}
  \FunctionTok{tm\_map}\NormalTok{(}\FunctionTok{content\_transformer}\NormalTok{(lemmatize\_words)) }\SpecialCharTok{\%\textgreater{}\%} \CommentTok{\# Lemmatization}
  \FunctionTok{tm\_map}\NormalTok{(stripWhitespace) }\SpecialCharTok{\%\textgreater{}\%} \CommentTok{\# Remove whitespace}
  \FunctionTok{DocumentTermMatrix}\NormalTok{() }\SpecialCharTok{\%\textgreater{}\%}
  \CommentTok{\# weightTfIdf() \%\textgreater{}\% \# Tf{-}Idf weighting (commented out)}
  \FunctionTok{removeSparseTerms}\NormalTok{(}\FloatTok{0.98}\NormalTok{) }\CommentTok{\# Adjust for sparsity}

\CommentTok{\# Check the dimensions of the DTM}
\FunctionTok{dim}\NormalTok{(dtm)}

\CommentTok{\# Convert the DTM to a data frame}
\NormalTok{dtm\_df }\OtherTok{\textless{}{-}} \FunctionTok{as.data.frame}\NormalTok{(}\FunctionTok{as.matrix}\NormalTok{(dtm))}

\CommentTok{\# Ensure dtm\_df and dt contain the same number of documents, then add Brand and Type}
\NormalTok{dtm\_df }\OtherTok{\textless{}{-}} \FunctionTok{bind\_cols}\NormalTok{(dtm\_df, dt[ , }\FunctionTok{c}\NormalTok{(}\StringTok{"Brand"}\NormalTok{, }\StringTok{"Type"}\NormalTok{)])}

\CommentTok{\# Split the data frame into training and testing sets and remove the Type column}
\NormalTok{dtm\_train }\OtherTok{\textless{}{-}} \FunctionTok{filter}\NormalTok{(dtm\_df, Type }\SpecialCharTok{==} \StringTok{"train"}\NormalTok{) }\SpecialCharTok{\%\textgreater{}\%} \FunctionTok{select}\NormalTok{(}\SpecialCharTok{{-}}\NormalTok{Type)}
\NormalTok{dtm\_test }\OtherTok{\textless{}{-}} \FunctionTok{filter}\NormalTok{(dtm\_df, Type }\SpecialCharTok{==} \StringTok{"test"}\NormalTok{) }\SpecialCharTok{\%\textgreater{}\%} \FunctionTok{select}\NormalTok{(}\SpecialCharTok{{-}}\NormalTok{Type)}

\CommentTok{\# Clean up the workspace}
\FunctionTok{rm}\NormalTok{(dtm, corpus, dtm\_df)}
\end{Highlighting}
\end{Shaded}

\subsection{Linear discriminant
analysis}\label{linear-discriminant-analysis}

Linear Discriminant Analysis (LDA) is a robust algorithm for
classifying, especially when distinguishing between more than
two classes. While logistic regression, which is another algorithm, is
widely used for binary classification (Komarek 2004), it faces
challenges with dealing with multiclass scenarios, where multinomial logistic
regression or LDA is more suitable (Torkkola 2001). In marketing,
both techniques are standard when the task is to describe market segments which have been derived previously (and which are built by using a clustering method); for an overview refer to Wedel and Kamakura
(2000) or Reutterer and Dan (2021). The word ``discriminate'' means to
separate and LDA uses a linear discrimination function to separate
classes, addressing some of the limitations of logistic regression, such
as instability in cases of well-separated classes or small datas ets. It
is important to distinguish LDA from another, non-related,
soft-clustering algorithm called latent Dirichlet allocation despite
sharing the same abbreviation. We will discuss the latter in the next
chapter. Classification LDA assumes that the variables follow a normal
distribution and that variances are equal across groups
(homoscedasticity), making outliers a potential issue, as with many
models (Balakrishnama and Ganapathiraju 1998).

Several algorithms can extend LDA's principles, including quadratic
discriminant analysis (Tharwat 2016), flexible discriminant analysis
(Hastie, Tibshirani, and Buja 1994), and regularized discriminant
analysis (Friedman 1989). In machine learning, model parameters are
estimated from training data to make predictions. The caret package in
R, standing for \textbf{c}lassification \textbf{a}nd \textbf{re}gression
\textbf{t}raining, facilitates this process with functions for data
splitting, pre-processing, feature selection, model tuning through
cross-validation, resampling, and estimating the importance of variables (Kuhn et al. 2013). For model tuning, caret's \texttt{trainControl()}
function is versatile, allowing the use of various validation techniques,
including K-fold cross-validation, a method we explore within the context of
LDA. Cross-validation divides the training set into K segments, or
folds, typically between 5 and 10, to train and validate the
model iteratively, enhancing its reliability and performance (Forman and Scholz
2010).

\begin{figure}

\centering{

\includegraphics[width=0.99\textwidth,height=\textheight]{figure/chapter6/cv.png}

}

\caption{\label{fig-cv}Supervised Learning - Cross Validation
\newline Authors' own figure}

\end{figure}%

Figure~\ref{fig-cv} shows the cross-validation process using a
5-fold example. This means there are five iterations, each incorporating
a unique combination of folds. In every iteration, a model is trained on
the training set and then applied to the validation set to make
predictions. The discrepancy between the model's predictions
and the ground-truth labels of the validation fold is measured as
mis-classification. The overall error rate \textbf{E} is calculated as
the average of the errors from each iteration \emph{i}. A common
misunderstanding is the belief that the best-performing model from these
iterations is selected for future classification tasks. However, this
approach does not fully align with best practice. Cross-validation
serves as a method to estimate how well the algorithm is likely to
perform on unseen data. It provides a measure of the model's
generalizability rather than of selecting a specific model for use. It is
also important to note that the test set remains untouched during
cross-validation and is reserved exclusively for final evaluation.

The final model, intended for operational use, is trained on the entire
training data set without folding, using the hyper-parameters chosen (or validated) during cross-validation. Once this model has been developed, it can
then be evaluated using the test set, with performance metrics providing
insights into its effectiveness with the help of a confusion matrix. A
confusion matrix is shown in tabular form and used to evaluate the
performance of classification models. It categorizes predictions into
four categories: true positives (correct positive predictions), false
positives (incorrectly predicted as positive), false negatives
(incorrectly predicted as negative), and true negatives (correct
negative predictions). When the class distribution is imbalanced, it is common to improve the overall accuracy with metrics such as precision, recall, and an F1-score.

To understand better how cross-validation works, let us link
image Figure~\ref{fig-cv} with table Table~\ref{tbl-cvtable} to pass
from the abstract concept from the image to a simple application in the from of a
table. We will use an example of a data set containing 10 documents. Before folding, the documents are usually shuffled at random to avoid any order-related bias.

\begin{longtable}[]{@{}
  >{\raggedright\arraybackslash}p{(\columnwidth - 4\tabcolsep) * \real{0.1067}}
  >{\raggedright\arraybackslash}p{(\columnwidth - 4\tabcolsep) * \real{0.2400}}
  >{\raggedright\arraybackslash}p{(\columnwidth - 4\tabcolsep) * \real{0.6533}}@{}}
\caption{5-Fold Cross-Validation with 10
Documents}\label{tbl-cvtable}\tabularnewline
\toprule\noalign{}
\begin{minipage}[b]{\linewidth}\raggedright
Fold
\end{minipage} & \begin{minipage}[b]{\linewidth}\raggedright
Validation Set
\end{minipage} & \begin{minipage}[b]{\linewidth}\raggedright
Training Set
\end{minipage} \\
\midrule\noalign{}
\endfirsthead
\toprule\noalign{}
\begin{minipage}[b]{\linewidth}\raggedright
Fold
\end{minipage} & \begin{minipage}[b]{\linewidth}\raggedright
Validation Set
\end{minipage} & \begin{minipage}[b]{\linewidth}\raggedright
Training Set
\end{minipage} \\
\midrule\noalign{}
\endhead
\bottomrule\noalign{}
\endlastfoot
Fold 1 & Doc1, Doc2 & Doc3, Doc4, Doc5, Doc6, Doc7, Doc8, Doc9, Doc10 \\
Fold 2 & Doc3, Doc4 & Doc1, Doc2, Doc5, Doc6, Doc7, Doc8, Doc9, Doc10 \\
Fold 3 & Doc5, Doc6 & Doc1, Doc2, Doc3, Doc4, Doc7, Doc8, Doc9, Doc10 \\
Fold 4 & Doc7, Doc8 & Doc1, Doc2, Doc3, Doc4, Doc5, Doc6, Doc9, Doc10 \\
Fold 5 & Doc9, Doc10 & Doc1, Doc2, Doc3, Doc4, Doc5, Doc6, Doc7, Doc8 \\
\end{longtable}

Table Table~\ref{tbl-cvtable} corresponds to the left portion of
Figure~\ref{fig-cv}, showing the 80\% of the data set designated as
the training set. In machine learning, ``training'' refers to the
process of teaching an algorithm to recognize patterns, relationships,
and features that lead to specific outcomes. For example, when applying
classification LDA to textual data, such as a document term matrix
(DTM), the predictors are the text features extracted from the reviews,
such as the frequencies of words or other linguistic properties. These features
help the algorithm identify which discriminate most between the two brands, revealing patterns that distinguish the brands in terms of
language, word usage, or specific terms.

This training process enables the development of models in \(K\) with
different configurations through \(K\)-fold cross-validation, where the
data set is divided into folds for sequential validation. Each fold
serves as a validation set while the remaining folds are used for
training. In our data set, cross-validation involves identifying the
conditions or features that distinguish between products A and B, as well as testing the model's ability to generalize.

But how effective is this model at predicting the classification of an
unseen review? For illustration, consider a simplified dataset comprising only 5 reviews. With \(K\)=5 , a 5-fold validation is effectively equivalent to leave-one-out cross-validation (LOOCV) in this case: It  uses each review sequentially as a validation point, initially, the first
review, then the second, and so on until the fifth. A unique model is
trained during each iteration, excluding the validation review from the
training process. The errors from each model are averaged to determine
the overall performance, with the algorithm optimizing itself to minimize this
error.

Subsequently, the model's ability to generalize is assessed by
evaluating its accuracy, or using other relevant metrics, for the test dataset. This was initially set aside, ensuring that it can provide a final,
unbiased evaluation of the model's performance.

\begin{Shaded}
\begin{Highlighting}[]
\CommentTok{\# Set 10{-}fold CV, save predictions, enable class probabilities}
\NormalTok{control }\OtherTok{\textless{}{-}} \FunctionTok{trainControl}\NormalTok{(}\AttributeTok{method =} \StringTok{"cv"}\NormalTok{, }\AttributeTok{number =} \DecValTok{10}\NormalTok{, }\AttributeTok{savePredictions =} \StringTok{"final"}\NormalTok{, }\AttributeTok{classProbs =} \ConstantTok{TRUE}\NormalTok{)}

\CommentTok{\# Set seed for reproducibility}
\FunctionTok{set.seed}\NormalTok{(seed)}

\CommentTok{\# Train LDA model, optimize for accuracy, 5 tuning iterations}
\NormalTok{lda }\OtherTok{\textless{}{-}} \FunctionTok{train}\NormalTok{(Brand }\SpecialCharTok{\textasciitilde{}}\NormalTok{ ., }\AttributeTok{data =}\NormalTok{ dtm\_train, }\AttributeTok{method =} \StringTok{"lda"}\NormalTok{, }\AttributeTok{trControl =}\NormalTok{ control, }\AttributeTok{metric=} \StringTok{"Accuracy"}\NormalTok{, }\AttributeTok{tuneLength =} \DecValTok{5}\NormalTok{)}

\CommentTok{\# Display model summary}
\NormalTok{lda}
\end{Highlighting}
\end{Shaded}

\begin{verbatim}
Linear Discriminant Analysis 

2939 samples
 804 predictor
   2 classes: 'Apple', 'Samsung' 

No pre-processing
Resampling: Cross-Validated (10 fold) 
Summary of sample sizes: 2645, 2646, 2645, 2645, 2645, 2645, ... 
Resampling results:

  Accuracy   Kappa    
  0.7727195  0.5429092
\end{verbatim}

The results from training the model can be summarized by entering the
variable name of where the model is stored. This summary includes the
number of classes, number of samples, predictor variables (the terms
from the document term matrix after the pre-processing), the type of
cross-validation used, and the evaluation metrics: accuracy and kappa.
Accuracy is particularly relevant for binary classification tasks,
especially if the data set is balanced (Koyejo et al. 2014), while kappa,
or Cohen's kappa, is more informative in scenarios where there is a class imbalance,
as it accounts for the possibility of the agreement occurring by chance
(Wang, Yang, and Xia 2019).

To apply the trained model to the test data set, we use the
\texttt{predict()} function. This not only classifies the new data but
can also provide probabilities for the assigning of each class, offering
deeper insights into the model's confidence in its predictions.

\begin{Shaded}
\begin{Highlighting}[]
\CommentTok{\# Predict with LDA model on test data}
\NormalTok{lda.predict }\OtherTok{\textless{}{-}} \FunctionTok{predict}\NormalTok{(lda, }\AttributeTok{newdata =}\NormalTok{ dtm\_test)}
\end{Highlighting}
\end{Shaded}

Now let us generate a confusion matrix by comparing the predictions made
by the classification model against the real, true labels of the data.
This matrix helps us to evaluate the model's performance by showing the
true positives, false positives, false negatives, and true negatives.
There are multiple metrics to assess how well a model fit can be
evaluated, through measures such as accuracy and error rate, as well as some
class-specific measures which evaluate its performance for the individual
classes rather than the overall data set, such as precision, recall, and
F-measures (Anandarajan et al. 2019).

The \texttt{caret} package in R provides comprehensive support for these
metrics for evaluating a model's performance thoroughly.

\begin{Shaded}
\begin{Highlighting}[]
\CommentTok{\# Evaluate model performance on the test set using a confusion matrix}
\NormalTok{cfm }\OtherTok{\textless{}{-}} \FunctionTok{confusionMatrix}\NormalTok{(lda.predict, }\FunctionTok{as.factor}\NormalTok{(dtm\_test}\SpecialCharTok{$}\NormalTok{Brand)) }
\FunctionTok{print}\NormalTok{(cfm)}
\end{Highlighting}
\end{Shaded}

\begin{verbatim}
Confusion Matrix and Statistics

          Reference
Prediction Apple Samsung
   Apple     325     108
   Samsung    47     257
                                          
               Accuracy : 0.7897          
                 95% CI : (0.7585, 0.8186)
    No Information Rate : 0.5047          
    P-Value [Acc > NIR] : < 2.2e-16       
                                          
                  Kappa : 0.5787          
                                          
 Mcnemar's Test P-Value : 1.441e-06       
                                          
            Sensitivity : 0.8737          
            Specificity : 0.7041          
         Pos Pred Value : 0.7506          
         Neg Pred Value : 0.8454          
             Prevalence : 0.5047          
         Detection Rate : 0.4410          
   Detection Prevalence : 0.5875          
      Balanced Accuracy : 0.7889          
                                          
       'Positive' Class : Apple           
                                          
\end{verbatim}

\begin{Shaded}
\begin{Highlighting}[]
\CommentTok{\# Obtain classification probabilities for each brand in the test set}
\NormalTok{lda.predict.p }\OtherTok{\textless{}{-}} \FunctionTok{predict}\NormalTok{(lda, }\AttributeTok{newdata =}\NormalTok{ dtm\_test, }\AttributeTok{type =} \StringTok{"prob"}\NormalTok{)}

\CommentTok{\# Display probabilities for the first six documents as percentages, rounded to two decimals}
\FunctionTok{head}\NormalTok{(}\FunctionTok{round}\NormalTok{(lda.predict.p }\SpecialCharTok{*} \DecValTok{100}\NormalTok{, }\DecValTok{2}\NormalTok{))}
\end{Highlighting}
\end{Shaded}

\begin{verbatim}
   Apple Samsung
4  90.13    9.87
5  92.70    7.30
8  99.99    0.01
11 67.03   32.97
16 93.08    6.92
20 98.90    1.10
\end{verbatim}

Examining the probabilities which a model uses to assign texts to specific
brands can provide deeper insights into the algorithm's confidence in
classifying objects. This analysis helps us understand that
inaccurate classification may arise from ambiguities and similarities in the
language used to describe different brands. In the following section, we
will explore a different classification algorithm and compare its
performance with linear discriminant analysis (LDA).

The accuracy of the test set and rounded to four decimals is
0.7897\% and quantifies the proportion of correct predictions of the
total predictions made by a model. Conversely, the error rate measures
the proportion of incorrect predictions, essentially serving as the
inverse of accuracy, calculated as Error rate = 1 - Accuracy. Here we
investigating metrics which assess how accurately individual classes are
predicted when evaluating class-specific performance. \textbf{Precision} measures
the proportion of correct predictions for a specific class out of all the
predictions made for that class, highlighting the model's accuracy in
identifying that class among all classes it attempts to identify.
\textbf{Recall}, on the other hand, focuses on the proportion of correct
predictions for a specific class out of all the actual instances of that
class, indicating the model's ability to capture all the instances.
The \textbf{F-measure}, or F1 score, harmonizes precision and recall into a
single metric, calculated as
\(F_{i}=2 \times \frac{precision\times recall}{precision+recall}\) where
\(i\) represents the class. This balance makes the F-measure a valuable
tool for assessing a model's overall effectiveness in class-specific
predictions, especially when using data sets which are imbalanced.

Using the confusion matrix, which assesses which
products were actual Apple or Samsung and were classified correctly (or
not), we can calculate all the measures mentioned above. See
Table~\ref{tbl-cfmtable} for better guidance.

\begin{longtable}[t]{>{}l|rr|>{}r}

\caption{\label{tbl-cfmtable}Confusion Matrix}

\tabularnewline

\\
\toprule
 & Actual Apple & Actual Samsung & Total Predicted\\
\midrule
\textbf{Predicted Apple} & 325 & 108 & 372\\
\textbf{Predicted Samsung} & 47 & 257 & 365\\
\midrule
\textbf{Total Actual} & 433 & 304 & 737\\
\bottomrule

\end{longtable}

With the numbers we obtained, let us calculate the accuracy:
\(\frac{325 + 257}{737}=0.7897\). The error rate is \(1-0.7897=0.2103\).
Precision for the Apple tablets is \(\frac{325}{372}=0,8737\) and
precision for the Samsung tablets is \(\frac{257}{365}=0.7041\). Recall
for the Apple products is \(\frac{325}{433}=0.7506\) while for the
Samsung products it is \(\frac{108}{304}=0.3553\). Calculating the
F-measure is straightforward by using the formula given in the paragraph above.

\subsubsection{A word about accuracy}\label{a-word-about-accuracy}

Assessing what constitutes ``good'' accuracy in machine learning is
nuanced, influenced by factors such as the specifics of the application,
the complexity of the problem, and class distribution within the data set. A
fundamental measure, baseline accuracy, represents the success rate of
the most naive model that always predicts the most frequent class. For a
model to be considered effective, its accuracy must substantially exceed
this baseline, indicating genuine learning rather than mere replication
of existing class distributions. Simple, well-defined problems might
warrant accuracies approaching perfection, whereas the messiness of
real-world data, characterized by noise and class overlap, means that for more complex issues accuracy will not meet the same level. The presence of imbalanced classes
further complicates matters, making high overall accuracy potentially
deceptive. Here, metrics such as precision, recall, and the F1 score,
particularly in relation to the performance of minority classes, offer a better
assessment of a model's efficacy. Moreover, the benchmark for ``good''
accuracy is not universal but varies across domains. For example, the
thresholds for acceptability in spam detection versus medical
diagnostics diverge significantly due to the difference in consequences resulting from
errors.

In the following section, we will explore a different classification
algorithm and compare its performance with linear discriminant analysis
(LDA).

\subsection{K-nearest neighbors (KNN)
classification}\label{k-nearest-neighbors-knn-classification}

The K-nearest neighbors (KNN) algorithm is simple and effective for tasks involving
classification and regression. Classification using KNN aims to
assign entities to predefined categories, whereas regression predicts a
numerical value. Our focus here will be on its use in
classification.

A unique element of KNN classficiation is that nothing is
ordained before the model is used. This means that, unlike algorithms
which learn parameters to make predictions for classification, KNN relies
on the data set itself, identifying the \(K\)'s most similar instances (or
``neighbors'') to a given point to make a prediction. New data are
evaluated during the running time against pre-existing data. The concept
of ``similarity'' is quantified through distance metrics, making KNN a
distance-based algorithm, and Euclidean distance is an excellent
candidate for this.

When the algorithm needs to make a prediction, it
searches through the training data set (i.e., the pre-existing data) to
find the instances (neighbors) nearest to \(K\).
The prediction is then made by assigning the new instance to one of the
predefined categories, based on the majority class among its nearest
neighbors.

Determining the optimal number of neighbors, \(K\), is essential for the
algorithm's performance, with the best value typically identified
through experimentation and validation techniques. KNN's principle is
straightforward: to classify an entity based on the majority vote of the
classes or average outcome of the target values of its \(K\)-nearest
neighbors (Brownlee 2016). This vote, or outcome, is based on the calculation of similarity between the new instance and its neighbors.
Note that in order to derive an unanimous decision of assignment, the
value for \(K\) must always be odd. With even values, the
algorithm would not know to which class the new element should be
assigned. While KNN is known for its ease of use and understanding, especially in scenarios with a small number of
predictors, its efficiency diminishes as the dimensions of the
data set increases, a phenomenon known as the ``curse of
dimensionality''. Further reading is recommended for a more detailed
exploration of KNN and its use (e.g., Tan 2006).

\begin{figure}[H]

{\centering \includegraphics[width=0.9\textwidth,height=\textheight]{figure/chapter6/knn.png}

}

\caption{KNN classification, \newline Authors' own figure}

\end{figure}%

In our example, we use exactly the same data set we prepared for
the LDA classification analysis. The training and testing data sets are
reused here, so the application of KNN classification is immediate. In this
analysis, we use the DTM to detect which textual features are
predictive of the two brands under study, focusing on how KNN classification identifies
patterns in the data.

\begin{Shaded}
\begin{Highlighting}[]
\CommentTok{\# Configure 10{-}fold CV with class probabilities}
\NormalTok{ctrl }\OtherTok{\textless{}{-}} \FunctionTok{trainControl}\NormalTok{(}\AttributeTok{method =} \StringTok{"cv"}\NormalTok{, }\AttributeTok{number =} \DecValTok{10}\NormalTok{, }\AttributeTok{classProbs =} \ConstantTok{TRUE}\NormalTok{)}

\CommentTok{\# Train kNN, optimizing accuracy, with 5 k values}
\FunctionTok{set.seed}\NormalTok{(seed)  }\CommentTok{\# Ensure reproducible results}
\NormalTok{knn }\OtherTok{\textless{}{-}} \FunctionTok{train}\NormalTok{(Brand }\SpecialCharTok{\textasciitilde{}}\NormalTok{ ., }\AttributeTok{data =} \FunctionTok{data.frame}\NormalTok{(dtm\_train), }\AttributeTok{method =} \StringTok{"knn"}\NormalTok{, }\AttributeTok{trControl =}\NormalTok{ ctrl, }\AttributeTok{metric =} \StringTok{"Accuracy"}\NormalTok{, }\AttributeTok{tuneLength =} \DecValTok{5}\NormalTok{)}

\CommentTok{\# Display model summary}
\NormalTok{knn}
\end{Highlighting}
\end{Shaded}

\begin{verbatim}
k-Nearest Neighbors 

2939 samples
 804 predictor
   2 classes: 'Apple', 'Samsung' 

No pre-processing
Resampling: Cross-Validated (10 fold) 
Summary of sample sizes: 2645, 2646, 2645, 2645, 2645, 2645, ... 
Resampling results across tuning parameters:

  k   Accuracy   Kappa    
   5  0.7985686  0.5916877
   7  0.7876831  0.5692389
   9  0.7778192  0.5488217
  11  0.7733881  0.5396807
  13  0.7638655  0.5198805

Accuracy was used to select the optimal model using the largest value.
The final value used for the model was k = 5.
\end{verbatim}

The system identified the optimal performance as \(k\)=5, where the
model's predictive accuracy reached 79.8569\%. It is important to note
that variations in the size and diversity of the data set could lead to improved
accuracy. Moreover, the pre-processing significantly influences
the results. In this analysis, employing a bag-of-words model (see
Section~\ref{sec-vsm-bow}) alongside lemmatization (see
Section~\ref{sec-reduction-of-inflectional-forms}) yielded the highest
accuracy. For a comprehensive comparison of the results under various conditions, see Table~\ref{tbl-knntable}. This discussion
presents the best-case scenario based on the current data set and
pre-processing methods. You are encouraged to experiment with different
pre-processing options to try and enhance the model's performance.

\begin{longtable}[]{@{}llll@{}}
\caption{KNN, various preprocessing
methods}\label{tbl-knntable}\tabularnewline
\toprule\noalign{}
k & Model & Word Root & Accuracy \\
\midrule\noalign{}
\endfirsthead
\toprule\noalign{}
k & Model & Word Root & Accuracy \\
\midrule\noalign{}
\endhead
\bottomrule\noalign{}
\endlastfoot
7 & BoW & Stem & 0.7842 \\
\textbf{5} & \textbf{BoW} & \textbf{Lem} & \textbf{0.7914} \\
9 & TF-IDF & Stem & 0.6189 \\
5 & TF-IDF & Lem & 0.6250 \\
\end{longtable}

We can also visualize the predictive performance with a varying \(K\)

\begin{Shaded}
\begin{Highlighting}[]
\CommentTok{\# plot and identify best k}
\FunctionTok{ggplot}\NormalTok{(knn)}
\end{Highlighting}
\end{Shaded}

\includegraphics{6.Chapter_files/figure-pdf/knn_plot-1.pdf}

The process for making predictions (or classifications) mirrors that of
the earlier example. To evaluate the algorithm's performance for the test
data set, we have used a confusion matrix.

\begin{Shaded}
\begin{Highlighting}[]
\CommentTok{\# Predict labels for the test dataset using the trained kNN model}
\NormalTok{knn.predict }\OtherTok{\textless{}{-}} \FunctionTok{predict}\NormalTok{(knn, }\AttributeTok{newdata =} \FunctionTok{data.frame}\NormalTok{(dtm\_test))}

\CommentTok{\# Evaluate the model\textquotesingle{}s performance on the test set using a confusion matrix}
\NormalTok{cfm }\OtherTok{\textless{}{-}} \FunctionTok{confusionMatrix}\NormalTok{(knn.predict, }\FunctionTok{as.factor}\NormalTok{(dtm\_test}\SpecialCharTok{$}\NormalTok{Brand))}
\end{Highlighting}
\end{Shaded}

In this case the accuracy for the testing set is 77.3406\%,
which is slightly worse than the LDA classification, but
otherwise more or less equivalent.

\subsection{Support-vector-machine
classification}\label{support-vector-machines-classification}

Support-vector-machine (SVM) classification is a well-known and widely used
procedure that performs well in binary classification and
multi-class classification (D. Cui and Curry 2005). It is
known to need little tuning when creating a model. We need to introduce
the concept of the maximal-margin classifier here to explain the internal workings
of SVM classification. Usually we have \(n\) input variables, such as the words in a
vocabulary. Classification means separation and is done through a
hyperplane. For example, in a 3-dimensional space, a hyperplane will
have 2 dimensions while in a 2-dimensional space the hyperplane has
1-dimension, basically a line. For simplicity let us take a two
dimensional space. In Figure~\ref{fig-svm}, we can see the optimal
solution SVM finds for the position of the separation line for
distinguishing two classes out of an infinity of possible solutions through which the
line can pass.

\begin{figure}

\centering{

\includegraphics[width=0.9\textwidth,height=\textheight]{figure/chapter6/svm.png}

}

\caption{\label{fig-svm}Support Vector Machines Classes Separation,
\newline Authors' own figure}

\end{figure}%

The model classifies new elements based on their position
relative to the decision boundary: Elements on one side of the line
(where the line-model equation yields a value greater than zero) are
assigned to one class, while those on the other side (with values less
than zero) are classified as part of the opposite class. Classification
becomes challenging for values that lie close to the decision boundary.
At the heart of this model is the distance between the decision boundary
and the nearest points on either side of it which belong to both classes. The
optimal decision boundary, known as the maximal-margin hyperplane, is
the one that maximizes this distance, ensuring the most precise
separation between (in our case: two) classes. The points nearest to
this hyperplane are called support vectors, as they ``support'' the
positioning of the hyperplane, playing an essential role in defining a
robust separator. During training, the model aims to identify this
optimal hyperplane, maximizing the margin for a given data set
(Brownlee 2016).

In scenarios where perfect separation is not feasible, the soft-margin
classifier offers a more flexible solution. It introduces a parameter,
denoted as C (for ``cost''), which adjusts the model's tolerance for
misclassifications. Tuning the C parameter allows for some points to
cross the boundary, effectively balancing the trade-off between a strict
separation and allowing for misclassifications. A higher C value tends to
reduce misclassifications by enforcing stricter margins.

The construction of the hyperplane for SVM is achieved through linear
algebra operations known as the dot product, also referred to as the
inner product. For instance, consider two vectors represented by {[}1,
2{]} and {[}3, 4{]}; their inner product is calculated as the sum of the
products of their corresponding elements: \(1\times 3 + 2\times 4\).
Prediction using SVM classification involves an equation that uses
the inner product between the input-feature values \((x)\) and the
support-vector values \((x_i)\),

\[f(x)=b_0 + \sum(w_i*(x,x_i))\]

while \(b_0\) and \(w_i\) are coefficients estimated from the model.

The prediction for a new data point involves using the dot product between the
input features and the support vectors. This operation, known as the
kernel function, essentially measures similarity, or from another
perspective, the distance between data points.

When dealing with a linear combination of input features, we employ a
linear kernel. However, to facilitate classification in scenarios where
the data are not separable by a line, inputs can be transformed into
higher dimensions using polynomial or radial kernel techniques. These
allow for complex separations with curved lines. We will explore the
performance of these various kernels to determine which offers the best
separation.

\subsubsection{Linear SVM.}\label{linear-svm.}

The linear SVM kernel is the simplest of various types of kernels
used in SVM models. It excels in scenarios with many features, making it
particularly well-suited for text classification. Gao
and Sun (2010) suggest that linear kernels are efficient and perform well
for high-dimensional data such as texts due to their straightforward
computation. Compared to more complex kernels, linear kernels are
faster, as they are less computationally intensive. The
formula for the linear kernel is:

\[K(x, x_i)=\sum(x*x_i)\] This equation represents the dot product of
the input vector \(x\) and a support vector \(x_i\), effectively
measuring the distance between new data points and the support vectors.
The training process for a model using a linear kernel typically
involves cross-validation, something you have seen before, to ensure
the model generalizes well for unseen data.

\begin{Shaded}
\begin{Highlighting}[]
\FunctionTok{library}\NormalTok{(kernlab)}
\CommentTok{\# Set 10{-}fold CV with class probabilities and save predictions}
\NormalTok{ctrl }\OtherTok{\textless{}{-}} \FunctionTok{trainControl}\NormalTok{(}\AttributeTok{method=}\StringTok{"cv"}\NormalTok{, }\AttributeTok{number=}\DecValTok{10}\NormalTok{, }\AttributeTok{savePredictions=}\StringTok{"final"}\NormalTok{, }\AttributeTok{classProbs=}\ConstantTok{TRUE}\NormalTok{)}

\CommentTok{\# Tuning grid for \textquotesingle{}C\textquotesingle{} parameter}
\NormalTok{tG }\OtherTok{\textless{}{-}} \FunctionTok{expand.grid}\NormalTok{(}\AttributeTok{C=}\FunctionTok{seq}\NormalTok{(}\DecValTok{0}\NormalTok{, }\DecValTok{2}\NormalTok{, }\AttributeTok{length=}\DecValTok{5}\NormalTok{))}

\CommentTok{\# Train SVM with linear kernel, tuning \textquotesingle{}C\textquotesingle{}}
\FunctionTok{set.seed}\NormalTok{(seed)}
\NormalTok{svm.linear }\OtherTok{\textless{}{-}} \FunctionTok{train}\NormalTok{(Brand}\SpecialCharTok{\textasciitilde{}}\NormalTok{., }\AttributeTok{data=}\NormalTok{dtm\_train, }\AttributeTok{method=}\StringTok{"svmLinear"}\NormalTok{, }\AttributeTok{trControl=}\NormalTok{ctrl, }\AttributeTok{preProcess =} \FunctionTok{c}\NormalTok{(}\StringTok{"center"}\NormalTok{,}\StringTok{"scale"}\NormalTok{), }\AttributeTok{tuneGrid=}\NormalTok{tG)}

\CommentTok{\# Predict and evaluate with linear kernel SVM}
\NormalTok{svm.linear.predict }\OtherTok{\textless{}{-}} \FunctionTok{predict}\NormalTok{(svm.linear, }\AttributeTok{newdata=}\NormalTok{dtm\_test)}
\FunctionTok{confusionMatrix}\NormalTok{(svm.linear.predict, }\FunctionTok{as.factor}\NormalTok{(dtm\_test}\SpecialCharTok{$}\NormalTok{Brand))}
\end{Highlighting}
\end{Shaded}

\begin{verbatim}
Confusion Matrix and Statistics

          Reference
Prediction Apple Samsung
   Apple     334      70
   Samsung    38     295
                                          
               Accuracy : 0.8535          
                 95% CI : (0.8258, 0.8782)
    No Information Rate : 0.5047          
    P-Value [Acc > NIR] : < 2.2e-16       
                                          
                  Kappa : 0.7067          
                                          
 Mcnemar's Test P-Value : 0.002855        
                                          
            Sensitivity : 0.8978          
            Specificity : 0.8082          
         Pos Pred Value : 0.8267          
         Neg Pred Value : 0.8859          
             Prevalence : 0.5047          
         Detection Rate : 0.4532          
   Detection Prevalence : 0.5482          
      Balanced Accuracy : 0.8530          
                                          
       'Positive' Class : Apple           
                                          
\end{verbatim}

Results can be significantly improved by adjusting
the parameters. However, this has a trade-off, meaning that more extensive tuning may lead to increased training times but can yield superior
outcomes. To assess the effectiveness of the tuned parameters, results
can be compared with those obtained using other models, such as LDA and
KNN, discussed earlier in this chapter. Metrics like accuracy,
precision, and the F1-score provide a basis for evaluating whether
tuning the parameters leads to meaningful improvements compared to other models.

\subsubsection{SVM radial}\label{svm-radial}

In the context of SVM, particularly when employing the radial-basis-function (RBF) kernel, two critical parameters are introduced: the cost
parameter (C) and the gamma parameter \(\gamma\) (sometimes referred to
as \(\sigma\)). The cost parameter C addresses the trade-off between
minimizing training errors and ensuring simplicity of the model for better
generalization, essentially controlling how much misclassification is
tolerated. On the other hand, the gamma parameter \(\gamma\) influences
the decision boundary's shape; lower values of \(\gamma\) may result in
more linear decision boundaries, while higher values allow for more
complex, curved boundaries that can closely fit the training data. The
influence of \(\gamma\) is captured by the RBF kernel formula:

\[K(x, x_i)=e^{-\gamma*\sum((x-x_i)^2)}\] This formula shows how
the similarity between any two data points \(x\) and \(x_i\) is
calculated, with \(\gamma\) modulating the impact of the distance
between these points on their calculated similarity. The values for
\(\gamma\) typically range from \(0<\gamma<1\) allowing for a wide
spectrum of decision-boundary complexity. The radial kernel, with its
ability to create complex polygons of separation, is a
powerful method, especially for data sets that are not
separable linearly, enabling the model to capture intricate patterns within the
data. The default value for \(\gamma\) is \(0.1\).

\begin{Shaded}
\begin{Highlighting}[]
\CommentTok{\# Set seed for reproducible results}
\FunctionTok{set.seed}\NormalTok{(seed)}

\CommentTok{\# Train SVM with radial basis function kernel on training data}
\NormalTok{svm.radial }\OtherTok{\textless{}{-}} \FunctionTok{train}\NormalTok{(Brand }\SpecialCharTok{\textasciitilde{}}\NormalTok{ ., }\AttributeTok{data =}\NormalTok{ dtm\_train, }\AttributeTok{method =} \StringTok{"svmRadial"}\NormalTok{, }\AttributeTok{trControl =}\NormalTok{ ctrl, }\AttributeTok{tuneLength =} \DecValTok{10}\NormalTok{)}

\CommentTok{\# Predict labels on test data}
\NormalTok{svm.radial.predict }\OtherTok{\textless{}{-}} \FunctionTok{predict}\NormalTok{(svm.radial, }\AttributeTok{newdata =}\NormalTok{ dtm\_test)}

\CommentTok{\# Evaluate predictions using a confusion matrix}
\FunctionTok{confusionMatrix}\NormalTok{(svm.radial.predict, }\FunctionTok{as.factor}\NormalTok{(dtm\_test}\SpecialCharTok{$}\NormalTok{Brand))}
\end{Highlighting}
\end{Shaded}

\begin{verbatim}
Confusion Matrix and Statistics

          Reference
Prediction Apple Samsung
   Apple     339      73
   Samsung    33     292
                                          
               Accuracy : 0.8562          
                 95% CI : (0.8287, 0.8807)
    No Information Rate : 0.5047          
    P-Value [Acc > NIR] : < 2.2e-16       
                                          
                  Kappa : 0.712           
                                          
 Mcnemar's Test P-Value : 0.0001519       
                                          
            Sensitivity : 0.9113          
            Specificity : 0.8000          
         Pos Pred Value : 0.8228          
         Neg Pred Value : 0.8985          
             Prevalence : 0.5047          
         Detection Rate : 0.4600          
   Detection Prevalence : 0.5590          
      Balanced Accuracy : 0.8556          
                                          
       'Positive' Class : Apple           
                                          
\end{verbatim}

The accuracy shows minimal difference when compared to the linear
case. Fine-tuning the parameters could potentially enhance the results
marginally. However, this increased flexibility comes at a cost: It may
lead to overfitting, where the model performs well using the training data
but struggles to generalize unseen data.

\subsubsection{Polynomial kernel SVM}\label{polynomial-kernel-svm}

The polynomial kernel extends the concept of the linear kernel, allowing
for more complex, non-linear decision boundaries. While it sometimes may
be less accurate or efficient than its linear counterpart, it
offers flexibility in capturing more complex patterns within the data.
The polynomial kernel is defined as:

\[K(x, x_i)=1+\sum(x*x_i)^d \] Here, \(d\) represents the degree of the
polynomial, which must be greater than 1 for the kernel to exhibit
non-linear properties. When \(d\) = 1, the polynomial kernel simplifies
to the linear kernel. The term `polynomial' suggests the ability to
model curved lines, enabling the separation of classes through curved
boundaries rather than straight lines. Obviously choosing different
values for \(d\) can lead to more complex decision boundaries and might
better fit the training data, but this will increase the risk of
overfitting. Also, a higher \(d\) will result in increased
computational complexity affecting training time and requiring more
computer resources.

\begin{Shaded}
\begin{Highlighting}[]
\CommentTok{\# Set seed for reproducible results}
\FunctionTok{set.seed}\NormalTok{(seed)}

\CommentTok{\# Train SVM with polynomial kernel on training data}
\NormalTok{svm.poly }\OtherTok{\textless{}{-}} \FunctionTok{train}\NormalTok{(Brand }\SpecialCharTok{\textasciitilde{}}\NormalTok{ ., }\AttributeTok{data =}\NormalTok{ dtm\_train, }\AttributeTok{method =} \StringTok{"svmPoly"}\NormalTok{, }\AttributeTok{trControl =}\NormalTok{ ctrl, }\AttributeTok{tuneLength =} \DecValTok{5}\NormalTok{)}

\CommentTok{\# Predict labels for test data using the trained model}
\NormalTok{svm.poly.predict }\OtherTok{\textless{}{-}} \FunctionTok{predict}\NormalTok{(svm.poly, }\AttributeTok{newdata =}\NormalTok{ dtm\_test)}

\CommentTok{\# Evaluate model performance with a confusion matrix}
\FunctionTok{confusionMatrix}\NormalTok{(svm.poly.predict, }\FunctionTok{as.factor}\NormalTok{(dtm\_test}\SpecialCharTok{$}\NormalTok{Brand))}
\end{Highlighting}
\end{Shaded}

\begin{verbatim}
Confusion Matrix and Statistics

          Reference
Prediction Apple Samsung
   Apple     335      71
   Samsung    37     294
                                          
               Accuracy : 0.8535          
                 95% CI : (0.8258, 0.8782)
    No Information Rate : 0.5047          
    P-Value [Acc > NIR] : < 2.2e-16       
                                          
                  Kappa : 0.7066          
                                          
 Mcnemar's Test P-Value : 0.001496        
                                          
            Sensitivity : 0.9005          
            Specificity : 0.8055          
         Pos Pred Value : 0.8251          
         Neg Pred Value : 0.8882          
             Prevalence : 0.5047          
         Detection Rate : 0.4545          
   Detection Prevalence : 0.5509          
      Balanced Accuracy : 0.8530          
                                          
       'Positive' Class : Apple           
                                          
\end{verbatim}

Presenting the three SVM options together helps us to make decisions about their use and
allows us to compare their relative performance quickly. Graphical
representation is usually a necessary aid in such cases. Presenting the
three SVM options together also helps us to decide which to
choose and allows for a quick comparison of their relative performance. For
instance, the polynomial kernel often performs well in scenarios with
non-linear relationships, but its effectiveness depends heavily on the
choice of degree and regularization parameters. Compared to the linear
kernel, it may offer more flexibility at the cost of increased
computation time and a higher risk of overfitting. The radial kernel can
model more complex patterns, but requires the careful tuning of parameters
like gamma to avoid underfitting or overfitting.

\begin{Shaded}
\begin{Highlighting}[]
\CommentTok{\# Aggregate resampling results for SVM models with different kernels}
\NormalTok{svm }\OtherTok{\textless{}{-}} \FunctionTok{resamples}\NormalTok{(}\FunctionTok{list}\NormalTok{(svm.linear, svm.radial, svm.poly))}

\CommentTok{\# Summarize the performance metrics of the models}
\FunctionTok{summary}\NormalTok{(svm)}
\end{Highlighting}
\end{Shaded}

\begin{verbatim}

Call:
summary.resamples(object = svm)

Models: Model1, Model2, Model3 
Number of resamples: 10 

Accuracy 
            Min.   1st Qu.    Median      Mean   3rd Qu.      Max. NA's
Model1 0.7823129 0.8010204 0.8231293 0.8192374 0.8366041 0.8469388    2
Model2 0.8061224 0.8154762 0.8282313 0.8240928 0.8320564 0.8367347    0
Model3 0.8265306 0.8375850 0.8534977 0.8550533 0.8656463 0.9013605    0

Kappa 
            Min.   1st Qu.    Median      Mean   3rd Qu.      Max. NA's
Model1 0.5625000 0.6011521 0.6451912 0.6370622 0.6716049 0.6921681    2
Model2 0.6097163 0.6291597 0.6551464 0.6465499 0.6627008 0.6718750    0
Model3 0.6506361 0.6739805 0.7061855 0.7088224 0.7305667 0.8019880    0
\end{verbatim}

\begin{Shaded}
\begin{Highlighting}[]
\CommentTok{\# Visualize the performance comparison using ggplot}
\FunctionTok{ggplot}\NormalTok{(svm)}
\end{Highlighting}
\end{Shaded}

\begin{figure}[H]

{\centering \includegraphics{6.Chapter_files/figure-pdf/SVMall-1.pdf}

}

\caption{Model performance comparison. \newline Authors' own figure}

\end{figure}%

Several observations emerge when comparing these models using
performance metrics such as accuracy, precision, recall, and the F1
scores. RBF and polynomial kernels often, but not consistently,
outperform the linear kernel in accuracy for non-linear data sets due to
their ability to deal with more complex patterns. This increased flexibility
can sometimes lead to a trade-off between precision and recall, with the
RBF kernel generally providing a good balance between the two for diverse
data sets. The F1 score, which combines precision and recall, tends to be
better for models using RBF and polynomial kernels in non-linear
settings, reflecting its greater adaptability.

These
observations are evident in the plot and output derived from the coding above. . For example, the RBF kernel shows improved F1
scores and a better balance of precision and recall compared to the
linear kernel. The polynomial kernel has a higher degree of accuracy in
scenarios where non-linear relationships dominate. The computational trade-offs and increased risk of overfitting
with more complex kernels, such as RBF and polynomial, are also
reflected in the results, requiring us to tune the performing
parameter and cross-validation for better performance.

The decision on which kernel to use should take into account the data set's characteristics, the computational resources available, and the balance required between interpreting the model and its complexity.
Experimentation and cross-validation remain essential for determining
the optimal kernel and parameter settings for any classification
task, ensuring the model is accurate and generalizable.

\subsection{Ensemble learning}\label{ensemble-learning}

What if you could put more than one model together to obtain better
prediction accuracy and higher overall performance? A technique that
can assist in such a situation is called ensemble learning. The
Oxford's dictionary gives a definition of the word ensemble as: ``a small group
of musicians, dancers or actors who perform together''. We can think of
our system as an orchestra. Each machine learning model is a musician,
and each does its best so that they all contribute to a
collective performance which is more than the sum of its parts.

Ensemble algorithms can be categorized into three types: bagging,
Bboosting, and stacking. Bagging involves training multiple models on
different sub-samples of a training data set, allowing for parallel
processing and reducing variance. Boosting, on the other hand,
trains models sequentially (in a chain), with each focusing on
correcting its predecessor's errors, effectively reducing bias. Unlike
the first two, stacking combines predictions from diverse models, using
an additional model (often called a meta-model) to determine the best
way to combine these predictions, in a similar to how an orchestra's
conductor ensures the musicians play together.

Each method has its pros and cons, and choosing between them
is determined by the challenge at hand. For further insights into
these methods, see Tattar (2018). Our focus here will be on
stacking. See the example in Figure~\ref{fig-ensemble}.

\begin{figure}

\centering{

\includegraphics[width=0.7\textwidth,height=\textheight]{figure/chapter6/ensemble.png}

}

\caption{\label{fig-ensemble}Stack ensemble algorithm (example),
\newline Authors' own figure}

\end{figure}%

Utilizing the \texttt{caret} package, we can assemble various models
under a meta-model that optimizes the combination of their predictions
for the best outcome. Following the setup of control attributes for
training, we then compile a list of machine-learning
algorithms, such as LDA, KNN, and linear SVM. Using the
\texttt{caretList()} function, these models are trained, and the accuracy of their predictions is assessed collectively, showing the power of
ensemble learning in achieving optimal results. This is accomplished by
specifying resampling methods, such as cross-validation, using the
\texttt{resamples()} function from the \texttt{caret} package. This allows
for comparison of the performance of the models and enhances the reliability of
the stacking.

\begin{Shaded}
\begin{Highlighting}[]
\CommentTok{\# Set a 10{-}fold cross{-}validation scheme, enabling class probability predictions and saving all predictions}
\NormalTok{ctrl }\OtherTok{\textless{}{-}} \FunctionTok{trainControl}\NormalTok{(}\AttributeTok{method=}\StringTok{"cv"}\NormalTok{, }\AttributeTok{number =} \DecValTok{10}\NormalTok{, }\AttributeTok{classProbs=}\ConstantTok{TRUE}\NormalTok{, }\AttributeTok{index=}\FunctionTok{createFolds}\NormalTok{(dtm\_train}\SpecialCharTok{$}\NormalTok{Brand, }\AttributeTok{k =} \DecValTok{10}\NormalTok{), }\AttributeTok{savePredictions =} \StringTok{"all"}\NormalTok{)}

\CommentTok{\# Define a list of machine learning algorithms for ensemble learning}
\NormalTok{algorithmList }\OtherTok{\textless{}{-}} \FunctionTok{c}\NormalTok{(}\StringTok{\textquotesingle{}lda\textquotesingle{}}\NormalTok{, }\StringTok{\textquotesingle{}knn\textquotesingle{}}\NormalTok{, }\StringTok{\textquotesingle{}svmLinear\textquotesingle{}}\NormalTok{)}

\FunctionTok{library}\NormalTok{(caretEnsemble)}
\FunctionTok{set.seed}\NormalTok{(seed)  }\CommentTok{\# Ensure reproducible results}

\CommentTok{\# Train models specified in the algorithm list with the defined control settings}
\NormalTok{models }\OtherTok{\textless{}{-}} \FunctionTok{caretList}\NormalTok{(Brand }\SpecialCharTok{\textasciitilde{}}\NormalTok{ ., }\AttributeTok{data =}\NormalTok{ dtm\_train, }\AttributeTok{trControl =}\NormalTok{ ctrl, }\AttributeTok{methodList =}\NormalTok{ algorithmList, }\AttributeTok{tuneLength =} \DecValTok{10}\NormalTok{)}

\CommentTok{\# Compare the models\textquotesingle{} performance using resampling}
\NormalTok{results }\OtherTok{\textless{}{-}} \FunctionTok{resamples}\NormalTok{(models)}

\CommentTok{\# Summarize the comparative performance of the models}
\FunctionTok{summary}\NormalTok{(results)}
\end{Highlighting}
\end{Shaded}

\begin{verbatim}

Call:
summary.resamples(object = results)

Models: lda, knn, svmLinear 
Number of resamples: 10 

Accuracy 
               Min.   1st Qu.    Median      Mean   3rd Qu.      Max. NA's
lda       0.5576560 0.5642479 0.6068053 0.6006904 0.6181474 0.6797732    1
knn       0.6378072 0.7121928 0.7251418 0.7181979 0.7377127 0.7682420    0
svmLinear 0.6748582 0.6875178 0.7009452 0.7135846 0.7219282 0.8370510    0

Kappa 
               Min.   1st Qu.    Median      Mean   3rd Qu.      Max. NA's
lda       0.1197228 0.1340173 0.2130066 0.2004211 0.2363048 0.3551089    1
knn       0.2567923 0.4135803 0.4418663 0.4262784 0.4670586 0.5308638    0
svmLinear 0.3434394 0.3676854 0.3986043 0.4233009 0.4407455 0.6735033    0
\end{verbatim}

The comparison of models can be represented by a graph. You can clearly
see the best model of the three. Be aware, that the lines
represent an interval, given by the accuracy of the 10 folds. The mean
is represented by the dot and it is the value to be used. Note that in our application, the KNN algorithm
performs slightly better than the SVM with a linear-kernel function, and
the accuracy is 0.7181979.

\begin{Shaded}
\begin{Highlighting}[]
\CommentTok{\# Plot the results}
\FunctionTok{ggplot}\NormalTok{(results)}
\end{Highlighting}
\end{Shaded}

\begin{figure}[H]

{\centering \includegraphics{6.Chapter_files/figure-pdf/dotplotensemble-1.pdf}

}

\caption{Algorithm perfomance comparison \newline Authors' own figure}

\end{figure}%

Now, used the models trained previously as the base layer, we
can employ a generalized linear model (GLM) as the meta-learner. The GLM
serves as a highly efficient prediction model, synthesizing the base
models' outputs to enhance the overall performance.

\begin{Shaded}
\begin{Highlighting}[]
\CommentTok{\# Configure 10{-}fold CV, saving predictions and probabilities}
\NormalTok{stackControl }\OtherTok{\textless{}{-}} \FunctionTok{trainControl}\NormalTok{(}\AttributeTok{method=}\StringTok{"cv"}\NormalTok{, }\AttributeTok{number=}\DecValTok{10}\NormalTok{, }\AttributeTok{savePredictions=}\ConstantTok{TRUE}\NormalTok{, }\AttributeTok{classProbs=}\ConstantTok{TRUE}\NormalTok{)}

\CommentTok{\# Stack models using GLM as meta{-}learner, optimizing for accuracy}
\NormalTok{stack.glm }\OtherTok{\textless{}{-}} \FunctionTok{caretStack}\NormalTok{(models, }\AttributeTok{method=}\StringTok{"glm"}\NormalTok{, }\AttributeTok{metric=}\StringTok{"Accuracy"}\NormalTok{, }\AttributeTok{trControl=}\NormalTok{stackControl, }\AttributeTok{tuneLength =} \DecValTok{10}\NormalTok{)}

\CommentTok{\# Display stacked model summary}
\NormalTok{stack.glm}
\end{Highlighting}
\end{Shaded}

\begin{verbatim}
A glm ensemble of 3 base models: lda, knn, svmLinear

Ensemble results:
Generalized Linear Model 

26451 samples
    3 predictor
    2 classes: 'Apple', 'Samsung' 

No pre-processing
Resampling: Cross-Validated (10 fold) 
Summary of sample sizes: 23806, 23806, 23806, 23806, 23805, 23806, ... 
Resampling results:

  Accuracy   Kappa    
  0.7640567  0.5256692
\end{verbatim}

The analysis demonstrates that the ensemble system, using
stacking, has enhanced the predictive accuracy, elevating it from
0.7181979 to 0.7640567. This improvement underscores the efficacy of
bringing together diverse models to achieve a superior performance.

Stacking, can lead to better accuracy, provided these models exhibit low
correlation in their predictions. The underlying rationale is that each model brings unique insights or captures different patterns
within the data, thereby complementing each other. When models generate
largely uncorrelated predictions, it indicates that they excel in
different elements of the task. This diversity enables the
meta-learner, or the classifier that integrates these predictions, to
discern the most effective way of amalgamating the strengths of each
model, thereby optimizing the overall predictive performance. However,
if the predictions from the sub-models are highly correlated (with
correlation coefficients exceeding 0.75), this implies that the models are
likely to be capturing similar patterns and making analogous predictions. In
such scenarios, the advantage of ensemble learning diminishes, as the
aggregation of predictions offers little incremental value over the
individual predictions. But this is not the case here, since the
correlations are modest. The incremental increase in accuracy to r
stack.glm\(error\) The accuracy, although modest, signifies the advantage of
employing a stacked ensemble over relying solely on a single SVM model.

\begin{Shaded}
\begin{Highlighting}[]
\CommentTok{\# Correlation matrix between model predictions}
\FunctionTok{modelCor}\NormalTok{(results)}
\end{Highlighting}
\end{Shaded}

\begin{verbatim}
          lda       knn svmLinear
lda         1        NA        NA
knn        NA 1.0000000 0.2160285
svmLinear  NA 0.2160285 1.0000000
\end{verbatim}

\begin{Shaded}
\begin{Highlighting}[]
\CommentTok{\# Scatterplot matrix (splom) to visually examine the correlations between model predictions}
\FunctionTok{splom}\NormalTok{(results)}
\end{Highlighting}
\end{Shaded}

\begin{figure}[H]

{\centering \includegraphics{6.Chapter_files/figure-pdf/correlationensempble-1.pdf}

}

\caption{Scatterplot matrix of model predictions}

\end{figure}%

The NA correlations between LDA and the other two models (knn and
svmLinear) suggest that the value was either not calculated, or is not
applicable. Although moderate, the positive correlation between knn and
svmLinear indicates some level of agreement in their predictions,
suggesting they might be recognizing similar patterns in the data.

\section{Neural networks}\label{neural-networks}

The fascination with understanding how our brain thinks and what can be done in computer
terms in order to mimic and outperform our reasoning has received a lot of attention in the academic literature. The
move from understanding the biological neural cells of the natural neural network
to the algorithmic constructs of an artificial neural network (ANN) is
remarkable. ANNs, a subset of machine-learning algorithms,
employ interconnected nodes across a number of layers to process data,
emulating the human brain's intricate network of neurons and synapses.
While these algorithms have proven incredibly potent, they are not
always the optimal solution for every problem. Their performance
can vary based on several factors, including the task's complexity, the
size of the data set, and the computational resources available. Despite these
limitations, ANNs have achieved remarkable success across a wide range of
applications, from image recognition (Quraishi, Choudhury, and De 2012),
to fraud detection (Krambia-Kapardis, Christodoulou, and Agathocleous 2010;
Asha and KR 2021), to medical diagnoses (Yen and Tsao 2024), to the
classification of natural-language patterns (Goldberg 2017), which is
our area of interest.

\begin{figure}[H]

{\centering \includegraphics[width=0.9\textwidth,height=\textheight]{figure/chapter6/neuron.png}

}

\caption{Neural networks}

\end{figure}%

In the following example, we will turn to the R package \texttt{nnet}.
\texttt{caret} can interface with it to train neural-network models,
allowing users to use \texttt{caret}'s convenient functions such as model
training and parameter tuning, while using the neural-network algorithms
provided by \texttt{nnet}. Despite the seemingly mystical aura
surrounding neural networks, the core steps in use them for
classification remain straightforward and familiar: data preparation,
setting up the model, network training, and performance evaluation.

\begin{Shaded}
\begin{Highlighting}[]
\CommentTok{\# Define a grid for tuning the decay parameter and the number of units in the hidden layer}
\NormalTok{nnetGrid }\OtherTok{\textless{}{-}} \FunctionTok{expand.grid}\NormalTok{(}\AttributeTok{decay =} \FunctionTok{c}\NormalTok{(}\DecValTok{0}\NormalTok{, }\FloatTok{0.01}\NormalTok{, .}\DecValTok{1}\NormalTok{), }\AttributeTok{size =} \FunctionTok{c}\NormalTok{(}\DecValTok{1}\SpecialCharTok{:}\DecValTok{10}\NormalTok{))}

\CommentTok{\# Determine the maximum number of units in the hidden layer from the grid}
\NormalTok{maxSize }\OtherTok{\textless{}{-}} \FunctionTok{max}\NormalTok{(nnetGrid}\SpecialCharTok{$}\NormalTok{size)}

\CommentTok{\# Calculate the maximum number of weights for the neural network}
\CommentTok{\# Formula: M(p+1)+M+1, where M is the number of hidden units, and p is the number of inputs}
\NormalTok{numWts }\OtherTok{\textless{}{-}} \DecValTok{1} \SpecialCharTok{*}\NormalTok{ (maxSize }\SpecialCharTok{*}\NormalTok{ (}\FunctionTok{length}\NormalTok{(dtm\_train) }\SpecialCharTok{+} \DecValTok{1}\NormalTok{) }\SpecialCharTok{+}\NormalTok{ maxSize }\SpecialCharTok{+} \DecValTok{1}\NormalTok{)}

\CommentTok{\# Set a 10{-}fold cross{-}validation scheme (see the ensemble method)}
\NormalTok{ctrl }\OtherTok{\textless{}{-}} \FunctionTok{trainControl}\NormalTok{(}\AttributeTok{method=}\StringTok{"cv"}\NormalTok{, }\AttributeTok{number =} \DecValTok{10}\NormalTok{, }\AttributeTok{classProbs=}\ConstantTok{TRUE}\NormalTok{, }\AttributeTok{index=}\FunctionTok{createFolds}\NormalTok{(dtm\_train}\SpecialCharTok{$}\NormalTok{Brand, }\AttributeTok{k =} \DecValTok{10}\NormalTok{), }\AttributeTok{savePredictions =} \StringTok{"all"}\NormalTok{)}

\CommentTok{\# Load the nnet package for neural network training}
\FunctionTok{library}\NormalTok{(nnet)}

\CommentTok{\# Set a random seed for reproducible results}
\FunctionTok{set.seed}\NormalTok{(seed)}

\CommentTok{\# Train the neural network model using the specified control parameters and tuning grid}
\NormalTok{NNModel }\OtherTok{\textless{}{-}} \FunctionTok{train}\NormalTok{(Brand }\SpecialCharTok{\textasciitilde{}}\NormalTok{ ., }\AttributeTok{data =}\NormalTok{ dtm\_train, }\AttributeTok{method =} \StringTok{"nnet"}\NormalTok{, }\AttributeTok{trace =} \ConstantTok{FALSE}\NormalTok{,}
                 \AttributeTok{trControl =}\NormalTok{ ctrl, }\AttributeTok{tuneGrid =}\NormalTok{ nnetGrid, }\CommentTok{\# Include the tuning grid in the training process}
                 \AttributeTok{na.action =}\NormalTok{ na.omit, }\AttributeTok{MaxNWts =}\NormalTok{ numWts)}

\CommentTok{\# Predict the test set using the trained neural network model}
\NormalTok{NNPredictions }\OtherTok{\textless{}{-}} \FunctionTok{predict}\NormalTok{(NNModel, dtm\_test)}

\CommentTok{\# Evaluate the model\textquotesingle{}s performance with a confusion matrix}
\FunctionTok{confusionMatrix}\NormalTok{(NNPredictions, }\FunctionTok{as.factor}\NormalTok{(dtm\_test}\SpecialCharTok{$}\NormalTok{Brand))}
\end{Highlighting}
\end{Shaded}

\begin{verbatim}
Confusion Matrix and Statistics

          Reference
Prediction Apple Samsung
   Apple     335      46
   Samsung    37     319
                                          
               Accuracy : 0.8874          
                 95% CI : (0.8623, 0.9093)
    No Information Rate : 0.5047          
    P-Value [Acc > NIR] : <2e-16          
                                          
                  Kappa : 0.7747          
                                          
 Mcnemar's Test P-Value : 0.3799          
                                          
            Sensitivity : 0.9005          
            Specificity : 0.8740          
         Pos Pred Value : 0.8793          
         Neg Pred Value : 0.8961          
             Prevalence : 0.5047          
         Detection Rate : 0.4545          
   Detection Prevalence : 0.5170          
      Balanced Accuracy : 0.8873          
                                          
       'Positive' Class : Apple           
                                          
\end{verbatim}

The accuracy achieved with this model surpasses previous results,
underscoring neural networks' potential for enhancing our ability to predict classifications. However, this improvement comes at a cost, notably in the
form of the significantly extended duration for training. Knowledge of this trade-off
between accuracy and computational efficiency is critical in when selecting a model.

\section{Fasttext}\label{fasttext}

FastText (Bhattacharjee 2018) is a powerful library developed by
Facebook's AI research team and designed for efficient text
classification and embedding words (more on this topic can be found in
Chapter~\ref{sec-extracting-information-from-embeddings}). During its
training process, it stands out by using sub-word information,
such as character n-grams, alongside whole words. This approach enables
FastText to handle texts with rare words or misspelling effectively,
capturing their meanings based on sub-word patterns. In R, the
\texttt{fastTextR} package serves as an interface with this library,
offering training for text classification models, generating
word embedding, and conducting nearest-neighbor searches within the
embedding space. FastText's efficiency is particularly notable when
working with large data sets, allowing for the rapid training of models without
compromising their accuracy. Another of FastText's key strengths is its
support for multiple languages, making it an excellent tool for
multilingual text analysis. FastText can be used for various NLP
tasks, including: sentiment analysis, where it helps identify emotional
tones in product reviews; language modeling, enhancing the performance
of predictive-text systems; and named-entity recognition, where it helps to classify the names of people, organizations, and locations in
unstructured text, which is an exciting feature for any marketing
analyst looking for brands or brand names in text. In the following
chapters, we will tackle all these topics. Meanwhile, here is an example
of how FastText can be used for classification tasks using our tablet
data set.

\begin{Shaded}
\begin{Highlighting}[]
\FunctionTok{library}\NormalTok{(fastTextR)}

\FunctionTok{set.seed}\NormalTok{(}\DecValTok{123}\NormalTok{)}
\CommentTok{\# Split data into training and testing sets}
\NormalTok{data\_train }\OtherTok{\textless{}{-}}\NormalTok{ dt[}\FunctionTok{sample}\NormalTok{(}\FunctionTok{nrow}\NormalTok{(dt), }\FunctionTok{nrow}\NormalTok{(dt) }\SpecialCharTok{*} \FloatTok{0.8}\NormalTok{), ]}
\NormalTok{data\_test }\OtherTok{\textless{}{-}}\NormalTok{ dt[}\SpecialCharTok{{-}}\FunctionTok{sample}\NormalTok{(}\FunctionTok{nrow}\NormalTok{(dt), }\FunctionTok{nrow}\NormalTok{(dt) }\SpecialCharTok{*} \FloatTok{0.8}\NormalTok{), ]}

\CommentTok{\# Prepare training and testing data}
\NormalTok{prepare\_data }\OtherTok{\textless{}{-}} \ControlFlowTok{function}\NormalTok{(data, file\_name) \{}
\NormalTok{  labels\_texts }\OtherTok{\textless{}{-}} \FunctionTok{paste0}\NormalTok{(}\StringTok{"\_\_label\_\_"}\NormalTok{, }\FunctionTok{tolower}\NormalTok{(data}\SpecialCharTok{$}\NormalTok{Brand), }\StringTok{" "}\NormalTok{, }\FunctionTok{tolower}\NormalTok{(data}\SpecialCharTok{$}\NormalTok{Content))}
\NormalTok{  normalized }\OtherTok{\textless{}{-}} \FunctionTok{ft\_normalize}\NormalTok{(labels\_texts)}
  \FunctionTok{writeLines}\NormalTok{(normalized, }\AttributeTok{con =}\NormalTok{ file\_name)}
\NormalTok{\}}

\FunctionTok{prepare\_data}\NormalTok{(data\_train, }\StringTok{"tablets.train"}\NormalTok{)}
\FunctionTok{prepare\_data}\NormalTok{(data\_test, }\StringTok{"tablets.test"}\NormalTok{)}

\CommentTok{\# Set training parameters}
\NormalTok{cntrl }\OtherTok{\textless{}{-}} \FunctionTok{ft\_control}\NormalTok{(}\AttributeTok{word\_vec\_size =} \DecValTok{10}\NormalTok{L, }\AttributeTok{learning\_rate =} \FloatTok{0.3}\NormalTok{, }\AttributeTok{max\_len\_ngram =} \DecValTok{2}\NormalTok{L, }
                    \AttributeTok{min\_count =} \DecValTok{1}\NormalTok{L, }\AttributeTok{nbuckets =} \DecValTok{10000000}\NormalTok{L, }\AttributeTok{epoch =} \DecValTok{10}\NormalTok{L, }\AttributeTok{nthreads =} \DecValTok{4}\NormalTok{L)}

\CommentTok{\# Train model and save}
\NormalTok{model }\OtherTok{\textless{}{-}} \FunctionTok{ft\_train}\NormalTok{(}\AttributeTok{file =} \StringTok{"tablets.train"}\NormalTok{, }\AttributeTok{method =} \StringTok{"supervised"}\NormalTok{, }\AttributeTok{control =}\NormalTok{ cntrl)}
\FunctionTok{ft\_save}\NormalTok{(model, }\StringTok{"tablets.bin"}\NormalTok{)}

\CommentTok{\# Load model and predict}
\NormalTok{model }\OtherTok{\textless{}{-}} \FunctionTok{ft\_load}\NormalTok{(}\StringTok{"tablets.bin"}\NormalTok{)}
\NormalTok{test\_pred }\OtherTok{\textless{}{-}} \FunctionTok{ft\_predict}\NormalTok{(model, }\AttributeTok{newdata =} \FunctionTok{ft\_normalize}\NormalTok{(}\FunctionTok{paste0}\NormalTok{(}\StringTok{"\_\_label\_\_"}\NormalTok{, }\FunctionTok{tolower}\NormalTok{(data\_test}\SpecialCharTok{$}\NormalTok{Content))), }\AttributeTok{k =} \DecValTok{1}\NormalTok{L)}

\CommentTok{\# Evaluate model}
\NormalTok{confusion\_matrix }\OtherTok{\textless{}{-}} \FunctionTok{table}\NormalTok{(}\AttributeTok{Truth =} \FunctionTok{paste0}\NormalTok{(}\StringTok{"\_\_label\_\_"}\NormalTok{, }\FunctionTok{tolower}\NormalTok{(data\_test}\SpecialCharTok{$}\NormalTok{Brand)), }
                          \AttributeTok{Prediction =}\NormalTok{ test\_pred}\SpecialCharTok{$}\NormalTok{label)}
\FunctionTok{print}\NormalTok{(confusion\_matrix)}
\end{Highlighting}
\end{Shaded}

\begin{verbatim}
                  Prediction
Truth              __label__apple __label__samsung
  __label__apple              381               25
  __label__samsung             30              300
\end{verbatim}

\begin{Shaded}
\begin{Highlighting}[]
\CommentTok{\# Calculate and print accuracy}
\NormalTok{accuracy }\OtherTok{\textless{}{-}} \FunctionTok{sum}\NormalTok{(}\FunctionTok{diag}\NormalTok{(confusion\_matrix)) }\SpecialCharTok{/} \FunctionTok{sum}\NormalTok{(confusion\_matrix)}
\FunctionTok{print}\NormalTok{(}\FunctionTok{sprintf}\NormalTok{(}\StringTok{"Accuracy: \%0.4f"}\NormalTok{, accuracy))}
\end{Highlighting}
\end{Shaded}

\begin{verbatim}
[1] "Accuracy: 0.9253"
\end{verbatim}

Compared to our previous results, the accuracy seems to be
exceptionally high, and this shows the model's power in classifying
documents quickly and efficiently.

There are other tools to classify text, such as Keras (Ramasubramanian
et al. 2019). It is a comprehensive neural-network library for both R
and Python. It is used extensively in building and training advanced
deep-learning models for text classification, supporting complex
architectures such as convolutional neural networks (CNNs) and recursive
neural networks (RNNs), although it can also be applied to textual data.

\section{Classification and its role in marketing
analytics}\label{classification-and-its-role-in-marketing-analytics}

Machine-learning classification techniques are useful in modern
marketing analytics, offering new tools to help understand and predict
consumer behavior (Ali Hakami and Hosni Mahmoud 2022). There is an
abundance of textual data generated by users from reviews, social media,
and surveys. Classifiers can understand patterns and trends that unveil
customer preferences and preceptions of brands. Through the algorithms presented
in this chapter, an analyst can categorize customer feedback, supporting
customer service and product development, see Noori (2021), and also classify sentiment in reviews. This topic will be tackled in depth
in Chapter~\ref{sec-sentiment-analysis}.

For example, methods like the ones we have described in this chapter can
categorize customer feedback, identifying common themes or pain points
associated with a product, paving the way for binary separations into
good/bad, or more nuanced understanding, such as classification using a Likert scale.
Another potential application could be in classifying the feedback for new
brands to understand which existing brands customers perceive as being similar.

\emph{Exercise starts}

\textbf{Exercise 1}: Naive Bayes Brand Classification. Use the tablets
data set to classify brands with a naive Bayes model. Evaluate the accuracy
and generate a confusion matrix.

\textbf{Exercise 2}: Tablet Categorization with Keras. Use a Keras
model to categorize tablets from their descriptions. Focus on the model
architecture and the evaluation of performance.

\textbf{Exercise 3}: FastText vs.~Naive Bayes. Select a dataset,
classify it with FastText, then compare its performance and training time
against a naive Bayes model.

\textbf{Exercise 4}: SVM Customer-Feedback Categorization. Use SVM to
categorize customer feedback. Pre-process the data, train the SVM, and assess
the accuracy.

\textbf{Exercise 5}: Ensemble Learning for Feedback Classification.
Use ensemble learning to classify feedback. Experiment with different
algorithms, noting the training time and accuracy. \emph{Exercise ends}

\textbf{Conclusion}

\emph{Conclusion starts} This chapter has explored various facets of text
classification using different machine-learning models and tools such as
LDA, KNN, SVM, ensemble-learning techniques, and FastText, as well as a few other tools. We used these models in R,
emphasizing the importance of pre-processing text, choosing the
right model architecture, and evaluating the performance of the model through
metrics such as accuracy, precision, recall, and the F1 score. The
discussion highlighted the nuanced considerations involved in determining
``good'' accuracy, factoring in baseline accuracy, problem complexity,
class distribution, domain-specific benchmarks, and the cost of
misclassifying. Specific exercises have applied these
concepts, ranging from brand classification and tablet categorization, to
comparing model efficiencies and execution times. \emph{Conclusion ends}

